\documentclass{article}

\usepackage[margin=1in]{geometry}
\usepackage{amsmath, amssymb}

\title{Working Document -- Conversational Advertisement}
\author{}
\date{}

\begin{document}

\maketitle

\section{Goal of this Document}

This document is a working space to:
\begin{itemize}
    \item Define research questions
    \item Formalize the problem mathematically
    \item Collect and cluster relevant papers
    \item Brainstorm possible experiments
\end{itemize}


\\ 

\fbox{
\parbox{0.9\linewidth}{
% PROBLEM SETTING

We study how conversational advertising interventions in LLM-based systems affect user experience and conversational dynamics in multi-turn interactions. Specifically, we analyze how ad format, insertion timing, and conversational context influence subjective experience, behavioral responses, and attention allocation during dialogue.
\\

Rather than optimizing advertising policies or platform-level objectives, our goal is to characterize the causal and predictive effects of advertising interventions on conversational systems. We focus on understanding how different ad designs influence trust, perceived intrusiveness, cognitive load, and overall user experience, as well as how they alter conversational trajectories over time.
\\

We further model user experience as a function of conversational state and intervention design, learning predictive mappings of the form:
\[
f(\text{AdType}, \text{Timing}, \text{User Personality}...) \rightarrow \text{User Experience}_t
\]
where user experience is operationalized through subjective (trust, intrusiveness, usefulness, overall satisfaction) and behavioral (conversion, abandonment, time to reply) outcomes.
\\

To complement behavioral and self-reported measures, we conduct a controlled user study combining EEG, eye-tracking, and personality profiling (OCEAN). This allows us to investigate whether neurophysiological signals provide additional explanatory power beyond behavioral and conversational features in modeling user experience under advertising interventions.
\\

Finally, we provide a structured experimental framework for studying conversational advertising in LLMs, enabling systematic analysis of how ad design, timing, and conversational context jointly shape user experience and interaction dynamics.

}
}
\newpage 


\section{Research Variables}.

\subsection{Indepdent Variables}

\begin{enumerate}
    \item \textbf{Ad Type (Format / Integration Mode)}  
    This is the primary experimental variable, capturing how advertising is presented within the conversational interface. Each type here represents a sensible instance of the taxonomy space we have previously defined with decreasing level of intrusiveness.
    
    \begin{enumerate}
        \item \textit{Inline persuasive suggestion}: promotional content subtly integrated within the assistant’s natural response flow.
        
        \item \textit{Sponsored conversational suggestion}: recommendation-style content embedded as part of conversational follow-ups (Perplexity-style).
        
        \item \textit{Sponsored recommendation}: clearly labeled suggested item inserted within the response (Bing-style integration).
        
        \item \textit{Explicit ad block}: highly salient, visually or structurally separated advertisement with a clear call-to-action (OpenAI-style placement).
    \end{enumerate}

    \item \textbf{Ad insertion timing} we will quantize this integer variable into 3 levels using data-driven arguments thanks to the works of \cite{}. We can use raw turns or normalized percentiles (p33, p66, p100...)

        \begin{enumerate}
            \item Early (1st turn)
            \item Mid (2-4 turns)
            \item Late (5+ turns)
        \end{enumerate}

    \item \textbf{Ad Greediness (frequency)} We will fix it to frequency=1 during the whole conversation. This variable becomes too noisy to study without a proper advertisement policy model. 
\end{enumerate}


\subsection{Dependent Variables}

We are going to define a hierarchy of dependent variables we are going to collect during the study, based on their importance and the amount of information they provide:

\begin{enumerate}
    \item Neurophysiological signals 
        \begin{itemize}
            \item EEG 1: Cognitive Load Index (CLI)
            \item EEG 2: Frontal Alpha Asymmetry (FAA)
            \item EEG 3: Frontal theta power (Workload marker, control)
            \item EEG 4: Engagement Index (control metric) 
            \item Visual 1: Fixation duration on the ad
            \item Visual 2: Time to first fixation 
            \item Visual 3: Number of gazes 
        \end{itemize}
        
    \item Behavioural signals 
        \begin{itemize}
            \item Ad click (Conversion)
            \item Time to reply/type 
            \item Abandonement 
        \end{itemize}
        
    \item Self-assessment signals
        It is crucial to understand that the semantics of the these signals refer to the "Trust" or "Usefulness" of the conversation not to the advertisement.
        \begin{itemize}
            \item Overall satisfaction (primary)
            \item Trust in system
            \item Perceived intrusiveness
            \item Perceived usefulness
        \end{itemize}

    \item (BONUS) Attention shift: change in conversational intent trajectory following ad insertion (e.g., transition between intent states after exposure to an ad)
\end{enumerate}

\subsection{Additional variables}

\begin{itemize}
    \item \textbf{User personality traits (OCEAN)}:     We include a reduced Big Five personality representation to capture stable user differences as an additional covariate.
\end{itemize}

\begin{itemize}
    \item \textbf{Task Type (BLOCKING / CONTEXTUAL FACTOR)}: Thanks to the user intent model (Thrad BERT) we can infer the intent of the conversation with 84\% accuracy, so this can tracked throughout the conversation $s_t$ and for the initial intent of the conversational tasks $s_0$.

    \item \textbf{Latent Task Advertising potential}: For some tasks covering an initial intent, there is a latent advertising potential score. In other words some intents (Searching places to go for trips) have more potential for ads than others (Debate about Israel tomato plantations).

\end{itemize}
    
\newpage
\section{Research Questions}

TLDR: This study investigates the effects of advertisements onto users, not advertisement policies!

% =========================================================
\subsection{1. Ad Type Effects on Advertising}
% =========================================================

\textbf{RQ1:} How does ad format influence subjective user experience (trust, perceived usefulness, and perceived intrusiveness)?

\textbf{RQ2:} How does ad format influence behavioural user experience (e.g., conversion and abandonment)?

\textbf{RQ3:} How does ad format influence conversational dynamics, measured as changes in conversational intent trajectory (Attention shifts)?

\textbf{RQ4:} How does conversational task type moderate the effect of ad format on subjective user experience?

\textbf{RQ5:} How do ad format and ad timing interact in shaping subjective user experience?

% =========================================================
\subsection{2. Temporal Effects on Advertisement}
% =========================================================

\textbf{RQ6:} How does ad timing (turn index) influence subjective user experience?

\textbf{RQ7:} How does ad timing influence behavioural user experience (e.g., conversion and abandonment)?

\textbf{RQ3:} How does ad timing influence conversational dynamics, measured as changes in conversational intent trajectory (Attention shifts)?

\textbf{RQ9:} Does the effect of ad timing on user experience vary across different ad formats?

% =========================================================
\subsection{3. User Heterogeneity}
% =========================================================

\textbf{RQ10:} Do personality traits (OCEAN) moderate the effect of ad format on subjective user experience?

\textbf{RQ11:} Do personality traits moderate the effect of ad timing on subjective user experience?

% =========================================================
\subsection{4. Neurophysiological Mechanisms}
% =========================================================

\textbf{RQ12:} Are ad format and timing associated with changes in cognitive load, attention, and affective arousal (EEG and eye-tracking measures)?

\textbf{RQ13:} Do neurophysiological signals provide additional explanatory power beyond self-reported and behavioural measures of user experience?

% =========================================================
\subsection{5. Predictive Analysis}
% =========================================================

\textbf{RQ14:} Can subjective user experience (trust, usefulness, intrusiveness) be accurately predicted from ad type, timing, and conversational context?

\textbf{RQ15:} Is ad timing (turn index heuristic) sufficient to approximate optimal ad placement in conversational systems?

\textbf{RQ16:} Do EEG and eye-tracking signals improve predictive performance of user experience models beyond behavioural and contextual features alone?

\newpage

\section{Contributions}
\label{Contributions}

\begin{enumerate}

\item \textbf{A Controlled Dataset for Conversational Advertising Interventions:}  
We introduce a structured dataset of multi-turn LLM interactions where ad type, insertion timing, and conversational context are systematically manipulated and paired with subjective, behavioral, and conversational outcome measures.

\vspace{0.2cm}

\item \textbf{A Multi-Output Model of User Experience under Conversational Ads:}  
We formulate and learn a supervised model that maps conversational state and ad interventions to subjective user experience (trust, usefulness, intrusiveness), behavioral outcomes (conversion, abandonment), and conversational dynamics.

\vspace{0.2cm}

\item \textbf{A Multimodal Analysis of Cognitive and Attentional Effects of Ad Interventions:}  
We integrate EEG and eye-tracking signals to study the neurophysiological correlates of conversational advertising and evaluate their incremental value beyond behavioral and self-reported measures.

\vspace{0.2cm}

\item \textbf{A Unified Experimental Framework for Studying Advertisement t in LLM Conversations:}  
We propose a controlled within-subject experimental design that enables systematic analysis of how ad format, timing, and conversational context jointly influence user experience and conversational dynamics. We also propose some advancements in the underlying theory defining a advertisement taxonomy and some novel metrics like attention shift.

\end{enumerate}

\newpage

\section{Experiments (WIP MUST BE REVISED)}

This work will be charactertized by only one experiment, we will have users experimenting with different engaging conversational tasks while ads while appear and the recordings of all these signals and self-reports will be logged for generating the final datasets.

\subsubsection{Experimental Design \& Session Structure}

Each participant interacts with an LLM across three independent conversational trials, each corresponding to a different task and ad condition. The study follows a within-subject design with counterbalanced ordering of ad types and tasks to mitigate ordering effects (Latin squares format kind of).
\\ 

Each experimental session is structured as follows with an approximated duration of 1.30h - 1.45h:

\begin{itemize}
    \item \textbf{Introduction and Consent (10 min)}
    \item \textbf{Sensor Setup, Calibration and Personality Assesment (15 min)}
    \item \textbf{Baseline Recording + Practice Trial (5 min)}
    \item \textbf{Main Experiment (3/4 trials, $\sim$15 min each)}
    \item \textbf{Final Survey and Debrief (5 min)}
\end{itemize}
\\ 

For each session we have multiple trials, all of which have the following standard structure:
\\ 

\begin{enumerate}
    \item \textbf{Task Presentation}: Participants are given an engaging task prompt designed to elicit natural conversation.

    \item \textbf{Conversational Interaction}: Participants engage in a conversation with the LLM for a minimum of 6 and a maximum of 10 turns. An advertisement is systematically introduced at fixed turns (turn 3 or 4 and turn 6/7) depending on the experimental condition.

    \item \textbf{Post-Trial Survey}: After each conversation, participants report perceived trust, intrusiveness, relevance, annoyance, and helpfulness using Likert-scale responses.
\end{enumerate}

\subsubsection{Task Design}

We design tasks to induce natural and engaging interactions across three broad intent genres:

\begin{itemize}
    \item \textbf{Informational}: e.g., improving a daily habit or learning how to optimize a personal workflow.
    
    \item \textbf{Transactional}: e.g., selecting a product under constraints or planning a purchase decision.
    
    \item \textbf{Reflective / Social}: e.g., discussing life decisions, goals, or personal dilemmas.
\end{itemize}

WORK IN PROGRESS TAKE SOME EXAMPLES TO CONSTRUCT THE FINAL 4 TASKS
\\ 
\begin{itemize}
    \item Find the best laptop you can find with a constrained budget of 500 dollars
    \item Organize a realistic trip this year 
    \item Find the perfect gift for someone special
    \item What is your biggest pain poin your day to day that an LLM could help you optimize? What area of your life you would wish to improve?
    \item What is the most annoying thing in your day to day? Interact with the LLM to gather ideas on how to fix it
    \item Find a new hobby that fits you
    \item Design your ideal weekend / routine
    \item What would happen if you did not go to work for 1 full day a week? (crazy what ifs to hook people)
    \item Where do you see yourself in 10 years? is it the place where you want to be? Otherwise what will you do to change it?
    \item What is an unpopular opinon of yours, try to debate it and win against an LLM.
    \item Ask for recommendations for any product you currently need (headphones, bag, phone...)
\end{itemize}

The final set of 4 task prompts we have concluded will elicit engagement from the subjects and belong to those intent genres are the following:

\begin{enumerate}
    \item 
\end{enumerate}


\subsubsection{Hypotheses \& Statistical Rigor}

Before going deep into experiment design one thing that must be adressed is the initial hypotheses we are testing against to make experiments more productive rather than "fishing expedition". To avoid this we will have to keep outcome variables low and plan beforehand:

\begin{itemize}
    \item Metrics
    \item Hypothesis
    \item Analysis Plan
\end{itemize}

Also with a expected study size of 50 people we must avoid multiple variables interaction study because that would lead to weak statistical results, thats why research question about user personality is secondary. Ad type is treated as the primary fixed effect, while task category and trial index are included as covariates.
\\
We define a single primary \textbf{independent variable (IV)}, Ad Type, capturing the format and integration of advertising within the conversational response. We consider two primary \textbf{dependent variables (DV)}:
\begin{itemize}
    \item \textbf{Trust}
    \item \textbf{Perceived Intrusiveness}
\end{itemize}

Secondary outcomes include:
\begin{itemize}
    \item Annoyance and perceived relevance
    \item Neurophysiological signals (EEG cognitive load, engagement)
    \item Eye-tracking measures of attention
\end{itemize}

We control for contextual \textbf{covariates} such as turn index, message and response length, conversation length, and inferred intent. User personality traits (OCEAN) and task type are treated as \textbf{moderating variables}, capturing heterogeneity in user sensitivity to different ad formats.
\\

Additionally, we collect neurophysiological signals (EEG and eye-tracking) to measure cognitive load and attention dynamics. These signals are used in exploratory analyses and predictive modeling.Some hypothesis testing similar as this one will be performed to aim to solve the research questions at hand:
\\

\textbf{H1:} More intrusive ad formats lead to a decrease in user trust.

\textbf{H2:} More intrusive ad formats lead to an increase in perceived intrusiveness.

\textbf{H3:} More intrusive ad formats reduce the probability of interaction continuation.

\textbf{H4:} Incorporating neurophysiological signals (EEG and eye-tracking) improves predictive performance of user response models compared to conversational features alone.

\textbf{H5:} User personality traits moderate the effect of ad type on user response (e.g., trust and perceived intrusiveness).

\subsubsection{Further design choices}

To mitigate ordering and carryover effects, both task order and ad type assignment are counterbalanced across participants. Task categories are rotated using a Latin-square design. Ad types are assigned such that each condition appears equally often in each trial position and across task categories.
\\

Moreover

\newpage

\section{Experiment Workflows (Summary)}

\subsection{Crowdsourcing Study}

\begin{enumerate}
    \item \textbf{Briefing and Introduction} \\
    Participants are introduced to the study, its purpose, and interaction guidelines.

    \item \textbf{LLM Test Conversation} \\
    Participants engage in a short practice interaction with the system to familiarize themselves with the interface.

    \item \textbf{Initial Self-Assessment} \\
    Baseline subjective measures are collected (e.g., trust, expectations, familiarity).

    \item \textbf{Personality Assessment (BFI-10)} \\
    Participants complete a short Big Five personality questionnaire.

    \item \textbf{Main Experimental Trials (4 Tasks)} \\
    For each participant, four conversational trials are conducted, each corresponding to a different task type. For each trial:
    \begin{itemize}
        \item Task briefing (short prompt describing the goal)
        \item Random assignment of experimental conditions:
        \[
        \text{AdType} \in \{4\}, \quad \text{Timing} \in \{3\}
        \]
        \item Conversational interaction with the LLM (multi-turn)
        \item Post-trial self-assessment (trust, intrusiveness, usefulness, satisfaction)
        \item Logging of behavioral metrics (e.g., response time, abandonment, interaction signals)
    \end{itemize}

    \item \textbf{Final Assessment} \\
    Overall user experience and satisfaction are collected.

    \item \textbf{Wrap-up and Offboarding} \\
    Participants are debriefed and exit the study.
\end{enumerate}

\newpage

\subsection{Laboratory User Study (EEG + Eye-Tracking)}

\begin{enumerate}
    \item \textbf{Briefing and Introduction} \\
    Participants receive an overview of the study and provide informed consent.

    \item \textbf{Sensor Setup and Personality Assessment} \\
    EEG and eye-tracking sensors are attached and calibrated. During this period, participants complete the BFI questionnaire (BFI-10 or BFI-44 depending on available time).

    \item \textbf{LLM Test Conversation and Calibration} \\
    Participants perform a short interaction with the system to ensure proper calibration of sensors and familiarity with the interface.

    \item \textbf{Initial Self-Assessment} \\
    Baseline subjective measures are collected.

    \item \textbf{Main Experimental Trials (4 Tasks)} \\
    The same trial structure as in the crowdsourcing study is followed. For each trial:
    \begin{itemize}
        \item Task briefing
        \item Random assignment of experimental conditions:
        \[
        \text{AdType} \in \{4\}, \quad \text{Timing} \in \{3\}
        \]
        \item Conversational interaction with the LLM
        \item Post-trial self-assessment
        \item Logging of behavioral metrics
        \item Recording of neurophysiological signals:
        \begin{itemize}
            \item EEG (cognitive load, engagement, etc.)
            \item Eye-tracking (fixations, gaze patterns)
        \end{itemize}
    \end{itemize}

    \item \textbf{Final Assessment} \\
    Participants report overall satisfaction and perceived experience.

    \item \textbf{Wrap-up, Sensor Removal, and Offboarding} \\
    Sensors are removed and participants are debriefed.
\end{enumerate}

\newpage


\section{Further work}

This section can be understood as a continuation of the original paper that builds on previous findings and using the same conversational advertisingtheoretical framework. Some ideas for future papers that can be researched:

\begin{itemize}
    \item Enhancing the dataset with data augmentation?
    \item RL Policy Learning (Ad insertion policies, greediness how many ads are placed constrained by trust...)
    \item Personalized Conversational Advertising under User Trait Heterogeneity
\end{itemize}

\newpage

\section{Implementation Details}

\begin{enumerate}
    \item We intentionally avoid optimized recommendation models to prevent confounding effects from learned ranking policies and instead rely on high-recall semantic retrieval to ensure relevance while preserving experimental control. This is why we will prefer neutral information retrieval rather than amazon CRS pretrained models, although their performance can be greater.

    \item GPU utilization: We need to fit a quantized LLM for inference (Qwen 3.6 30B) an embedding (Qwen 3 Embedding 8B) and reranker model (Qwen 3 Reranker 8B) for RAG, and a BERT/Distillbert for intent management, so system design must be managed wisely to avoid fragmentation of VRAM and increased latency times due to PCIe data transfers:
        \begin{itemize}
            \item GPU 0 (40GB): Quantized Qwen 3.6 35b a3b (<24gb) + KV-cache
            \item GPU 1 (40GB): Embedding Model (8B) + Reranker Model (8B) 
            \item CPU (112 nodes, 512 GB RAM): BERT + UI + Sensors management + Logging + VectorDB
        \end{itemize}
\end{enumerate}



\newpage

\begin{thebibliography}{9}

\bibitem{xu2026ad}
Shengwei Xu, Zhaohua Chen, Xiaotie Deng, Zhiyi Huang, and Grant Schoenebeck.
\textit{Ad Insertion in LLM-Generated Responses}.
arXiv preprint arXiv:2601.19435, 2026.
doi: 10.48550/arXiv.2601.19435.

\bibitem{heineking2026detecting}
Sebastian Heineking, Wilhelm Pertsch, Ines Zelch, Janek Bevendorff, Benno Stein, Matthias Hagen, and Martin Potthast.
\textit{Detecting RAG Advertisements Across Advertising Styles}.
arXiv preprint arXiv:2603.04925, 2026.
doi: 10.48550/arXiv.2603.04925.

\bibitem{zhao2025exploring}
Xiaoyan Zhao, Yang Deng, Wenjie Wang, Hongzhan Lin, Hong Cheng, Rui Zhang, See-Kiong Ng, and Tat-Seng Chua.
\textit{Exploring the Impact of Personality Traits on Conversational Recommender Systems: A Simulation with Large Language Models}.
arXiv preprint arXiv:2504.12313, 2025.
doi: 10.48550/arXiv.2504.12313.

\bibitem{feizi2023online}
Soheil Feizi, MohammadTaghi Hajiaghayi, Keivan Rezaei, and Suho Shin.
\textit{Online Advertisements with LLMs: Opportunities and Challenges}.
arXiv preprint arXiv:2311.07601, 2023.
doi: 10.48550/arXiv.2311.07601.

\bibitem{martina2025distillrecdial}
Alessandro Francesco Maria Martina, Alessandro Petruzzelli, Cataldo Musto, Marco de Gemmis, Pasquale Lops, and Giovanni Semeraro.
\textit{DistillRecDial: A Knowledge-Distilled Dataset Capturing User Diversity in Conversational Recommendation}.
In \textit{Proceedings of the 19th ACM Conference on Recommender Systems}, 2025.
ACM.
doi: 10.1145/3705328.3748161.

\bibitem{openai_ads_chatgpt}
OpenAI. \emph{Ads in ChatGPT}. OpenAI Help Center. Available at:
\texttt{https://help.openai.com/en/articles/20001047-ads-in-chatgpt}.
Accessed: 2026-04-19.

\bibitem{rainie2025llmusers}
L. Rainie.
\newblock \textit{Close Encounters of the AI Kind: The Increasingly Human-Like Way People Are Engaging with Language Models}.
\newblock Imagining the Digital Future Center, Elon University, 2025.
\newblock Available online.

\bibitem{meguellati2025llmads}
Elyas Meguellati, Stefano Civelli, Lei Han, Abraham Bernstein, 
Shazia Sadiq, and Gianluca Demartini.
\newblock LLM-Generated Ads: From Personalization Parity to Persuasion Superiority.
\newblock \emph{arXiv preprint arXiv:2512.03373}, 2025.

\bibitem{borlund2010simulatedworktask}
Pia Borlund and Jesper W. Schneider.
\newblock Reconsideration of the simulated work task situation: a context instrument for evaluation of information retrieval interaction.
\newblock \textit{Proceedings of the Third Symposium on Information Interaction in Context (IIiX '10)}, 2010.
\newblock ACM.
\newblock doi: 10.1145/1840784.1840808.

\bibitem{kosmyna2025brainchatgpt}
Nataliya Kosmyna, Eugene Hauptmann, Ye Tong Yuan, Jessica Situ,
Xian-Hao Liao, Ashly Vivian Beresnitzky, Iris Braunstein, and Pattie Maes.
\newblock \textit{Your Brain on ChatGPT: Accumulation of Cognitive Debt when Using an AI Assistant for Essay Writing Task}.
\newblock \emph{arXiv preprint arXiv:2506.08872}, 2025.

\bibitem{baradari2025neurochat}
Dünya Baradari, Nataliya Kosmyna, Oscar Petrov,
Rebecah Kaplun, and Pattie Maes.
\newblock NeuroChat: A Neuroadaptive AI Chatbot for Customizing Learning Experiences.
\newblock \emph{arXiv preprint arXiv:2503.07599}, 2025.

\bibitem{saffaryazdi2025empathetic}
Nastaran Saffaryazdi, Tamil Selvan Gunasekaran,
Kate Laveys, Elizabeth Broadbent, and Mark Billinghurst.
\newblock Empathetic Conversational Agents: Utilizing Neural and Physiological Signals for Enhanced Empathetic Interactions.
\newblock \emph{International Journal of Human--Computer Interaction}, 42(6), 2026.
\newblock doi: 10.1080/10447318.2025.2540500.
\end{thebibliography}

\end{document}